\documentclass[10pt]{article}
\usepackage[T1]{fontenc}
\usepackage[utf8]{inputenc}
\usepackage[french]{babel}
\usepackage{amsmath,amssymb,mathrsfs,mathpazo}
\usepackage{geometry}
\geometry{a4paper,margin=2cm}
\usepackage{graphicx,array,tabularx,tabularray,colortbl,multirow,multicol,booktabs,enumitem,pifont,ulem,fancybox,diagbox,cancel}
\usepackage{pstricks,pstricks-add,pst-plot,pst-tree,pst-node,pst-text,pst-func,pst-math,pst-bspline,pst-3dplot,pst-eucl,pst-solides3d}
\let\pscalculate\undefined
\usepackage{pst-bezier}
\usepackage{tikz,pgfplots,tkz-tab}
\usetikzlibrary{arrows,arrows.meta,patterns,calc,positioning,trees,decorations.markings,shapes.geometric,shapes.misc,shapes.arrows,intersections,decorations.pathmorphing,decorations.pathreplacing,backgrounds,shadows,babel}
\pgfplotsset{compat=1.18}
\usepgfplotslibrary{fillbetween}
\usepackage[np]{numprint}
\usepackage{listings,makecell,pas-tableur,twoopt}
\let\barre\undefined
\usepackage[pstricks]{tablvar}
\definecolor{darkgreen}{RGB}{0,100,0}
\providecommand{\red}{\color{red}}
\providecommand{\blue}{\color{blue}}
\providecommand{\green}{\color{green}}
\providecommand{\black}{\color{black}}
\providecommand{\white}{\color{white}}
\providecommand{\gray}{\color{gray}}
\providecommand{\cyan}{\color{cyan}}
\providecommand{\magenta}{\color{magenta}}
\providecommand{\yellow}{\color{yellow}}
\providecommand{\violet}{\color{violet}}
\providecommand{\orange}{\color{orange}}
\providecommand{\euro}{\texteuro}
\providecommand{\EUR}{\texteuro}
\providecommand{\ieme}{\textsuperscript{e}}
\providecommand{\degres}{\ensuremath{^\circ}}
\providecommand{\degree}{\ensuremath{^\circ}}
\providecommand{\vv}[1]{\overrightarrow{#1}}
\providecommand{\Coord}[1]{#1}
\providecommand{\decosix}{}
\providecommand{\starredbullet}{\ensuremath{\bullet}}
\pagestyle{empty}
\setlength{\parindent}{0pt}
\usepackage{adjustbox}\DeclareUnicodeCharacter{2236}{:}\DeclareUnicodeCharacter{2190}{\ensuremath{\leftarrow}}
\providecommand{\R}{\mathbb{R}}
\providecommand{\N}{\mathbb{N}}
\providecommand{\Z}{\mathbb{Z}}
\providecommand{\Q}{\mathbb{Q}}
\providecommand{\D}{\mathbb{D}}
\providecommand{\C}{\mathbb{C}}
\providecommand{\e}{{\mathrm{e}}}
\providecommand{\vect}[1]{\overrightarrow{#1}}
\providecommand{\vectt}[1]{\overrightarrow{\mathrm{#1}}}
\providecommand{\vvt}[1]{\overrightarrow{\mathrm{#1}}}
\providecommand{\np}[2][]{{#2~#1}}
\providecommand{\vv}[1]{{\overrightarrow{#1}}}
\renewcommand{\d}{{\mathrm{d}}}
\providecommand{\degres}{{^\circ}}
\providecommand{\degre}{{^\circ}}
\providecommand{\blue}{\color{blue}}
\providecommand{\red}{\color{red}}
\providecommand{\enskip}{\quad}
\providecommand{\llbracket}{\lbrack\!\lbrack}
\providecommand{\rrbracket}{\rbrack\!\rbrack}
\providecommand{\euro}{{\text{€}}}
\providecommand{\no}{{\text{n°}}}
\providecommand{\interval}[2]{{[#1;#2]}}
\providecommand{\intervFO}[2]{{[#1;#2[}}
\providecommand{\dotfill}{\ldots}

\providecommand{\ds}{\displaystyle}
\providecommand{\textding}[1]{#1}
\providecommand{\Oij}{\(\mathrm{O};\vec{\imath},\vec{\jmath}\)}
\providecommand{\Oijk}{\(\mathrm{O};\vec{\imath},\vec{\jmath},\vec{k}\)}
\providecommand{\Ouv}{\(\mathrm{O};\vec{u},\vec{v}\)}
\providecommand{\intervFF}[2]{\ensuremath{[#1;#2]}}
\renewcommand{\euro}{\ifmmode\text{\texteuro}\else\texteuro\fi}
\DeclareUnicodeCharacter{00B0}{\ifmmode^{\circ}\else\textdegree\fi}
\DeclareUnicodeCharacter{20AC}{\ifmmode\text{\texteuro}\else\texteuro\fi}
% Énoncés : Éducation nationale. Transcription et corrigés : APMEP, Denis Vergès et contributeurs. Découpage et classement : CoopMaths, Sébastien Lozano, Éric Elter et contributeurs. Adaptation de la présentation : Maths Explorer.
% CC BY-SA : https://coopmaths.fr/www/about/
% https://coopmaths.fr/alea/static/bac/2022/tex/bac_2022_05_sujet1_etrangers_4_cor.tex
\begin{document}
\noindent\textit{APMEP / CoopMaths. CC BY-SA. Présentation adaptée par Maths Explorer.}\par\medskip
\bigskip









\textbf{Partie A}

\medskip

















\medskip

\begin{enumerate}
\item Recopier et compléter le tableau suivant avec les probabilités correspondantes.

\begin{center}
\begin{tabularx}{0.45\linewidth}{|*{4}{>{\centering \arraybackslash}X|}}\hline
				&$A$ 	&$\overline{A}$	&Total\\ \hline
\black $B$				&\blue 0,05		&\blue 0,15				&\black 0,2\\ \hline
\black $\overline{B}$	&\blue 0,05		&\black 0,75				&\blue 0,8\\ \hline
\black Total			&\black 0,1		&\blue 0,9				&\black 1\\ \hline
\end{tabularx}
\end{center}

\item
	\begin{enumerate}
		\item 
La probabilité qu'une paire de verres, prélevée au hasard dans la production, présente un défaut pour au moins un des deux traitements T1 ou T2 est l'évènement contraire de l'évènement \og une paire de verres ne présente aucun des deux défauts \fg, donc sa probabilité est égale $1 - 0,75 = 0,25$. Ou encore $P(A \cup B) = 1 - \left(\overline{A} \cap \overline{B} \right)$.
		\item 
		
$\bullet~~$On a $P\left(\overline{A}\right) = 1 - P(A) = 1 - 0,1 = 0,9$ ;

$\bullet~~$On a $P\left(\overline{B}\right) = 1 - P(B) = 1 - 0,2 = 0,8$.

Donc $P\left(A \cap \overline{B}\right) = 0,8 - 0,75 = 0,05$ et 

$P\left(B \cap \overline{A}\right) = 0,9 - 0,75 = 0,15$.

Donc la probabilité qu'une paire de verres, prélevée au hasard dans la production, présente deux  défauts  est égale à :

$P(A \cap B) = 1 - P\left(A \cap \overline{B}\right)  -  P\left(B \cap \overline{A}\right) - P\left(\overline{A} \cap \overline{B} \right) = 1 - 0,05 - 0,15  - 0,75 = 0,05$.
		\item 
$P(A) \times P(B) = 0,1 \times 0,2 = 0,02$ et $P(A \cap B) = 0,05$.

On a donc $P(A \cap B) \ne P(A) \times P(B)$ : les évènements $A$ et $B$ ne sont pas indépendants.
	\end{enumerate}	
\item 
On a $P\left(A \cap \overline{B}\right) + P\left(B \cap \overline{A}\right)= 0,05 + 0,15 = 0,2$
\item 
On a $P_{A}\left(B \right) = \dfrac{P(A \cap B)}{P(A)} = \dfrac{0,05}{0,1} = \dfrac12$.
\end{enumerate}

\bigskip

\textbf{Partie B}

\medskip







\begin{enumerate}
\item 
La production est suffisamment importante pour que la probabilité d'avoir le déf	ut $T1$ est égale à 0,1. Comme il y a 50 tirages indépendants, la variable aléatoire $X$ suit une loi binomiale de paramètres $n = 50$ et $p = 0,1$.
\item 
On sait que $P(X = 10) = \binom{50}{10}\times 0,1^{10} \times (1 - 0,1)^{50 - 10}$.

La calculatrice donne $P(X = 10)  \approx 0,015$ au millième près.


\item 
La moyenne est l'espérance de la variable $X$ et on sait que $E(X) = n \times p = 50 \times 0,1 = 5$.
\end{enumerate}

\end{document}
